\cotefile{156-8}
\scanpage{114}
\mspage{105}
% archive sheet number: 114 (upper left)
% L1
Je vais prouver
\[
\begin{aligned}
Z<X&\Longrightarrow Z\in\operatorname{supp}^{\circ}(\Phi)\\
Z>X&\Longrightarrow Z\in\operatorname{supp}^{\circ}(\Psi).
\end{aligned}
\]
% L4-6
\deleted{Pour ceci,} Prouvons p. ex. le premier, $Z<\unsure{x|X}$ (et \inserted{\unsure{ou simplement}} $Z\in\operatorname{supp}^{\circ}(\Phi\cup\Psi)$)
Soit donc
\[
Z'\in\operatorname{cosupp}^{\circ}(\Phi),\quad\text{i.e.}\quad Z'\mathrel{\mid\circ\mid}\Phi
\]
% L7-10
nous \deleted{\unsure{allons}} devons prouver $Z\mathrel{\unsure{\rangle\circ\mid|\mid\circ\mid}}Z'$. Nous \deleted{Soit} savons que c'est gagné si $Z'\mathrel{\mid\circ\mid}\Phi\cup\Psi$ \inserted{\unsure{resp.} $Z\mathrel{\mid\circ\mid}\Psi$ \unsure{aussi}} (et non seulement $Z'\mathrel{\mid\circ\mid}\Phi$).

% L11-15
\deleted{Mais soit}

\unsure{Quitte à|Cas A} : ``\unsure{agrandir}'' $X$, \unsure{on}
\[
\delta X=\{\underbrace{\operatorname{ex}\delta\Phi}_{a},\underbrace{\operatorname{or}\delta\Psi}_{b}\}
\]
i.e.
\[
\left\{\begin{aligned}
\deleted{\unsure{S}}X&=P_a&&\text{si}\quad a=b\\
X&=S_{ab}&&\text{si}\quad a<b
\end{aligned}\right.
\]
% L16-27
\deleted{\illegible{1}} Comme $Z'\mathrel{\mid\circ\mid}\Phi$, il s'ensuit aisément que $Z'\mathrel{\mid\circ\mid}P_a$ \uline{ou} $Z'=P_a$ \unsure{.}

% slanted
\margin{\uline{Lemme} \inserted{(\unsure{trivial})} $Z'\mathrel{\mid\circ\mid}Y$, et $\deleted{\unsure{P}}P\mathrel{\unsure{\triangleleft|\unlhd}}Y$ alors $Z'\mathrel{\unsure{\parallel|\mid\circ\mid}}P$ \uline{ou} $Z'=\deleted{\unsure{\{}}P\deleted{\unsure{\}}}$}

(Car on aura $\delta\Phi\cup\delta Z'\in\operatorname{Drap}(\unsure{\mathcal{L}|L})$, a fortiori $\delta\deleted{\unsure{\Phi}}Z'\cup\{P_a\}\in\operatorname{Drap}\mathcal{L}$, \deleted{\unsure{si on}} Si on avait $Z'=\deleted{\unsure{\{}}P_a$, \unsure{aurions|avions} par hypothèse $P_a\unsure{\ne|=|\in}Z'\mathrel{\mid\circ\mid}\Phi$, et $P_a\mathrel{\mid\circ\mid}\Psi$ puisque $X<\Psi$, on a $Z'\mathrel{\mid\circ\mid}\Phi\cup\Psi$, et on gagne. Sinon On aura $Z'<_{\mathrm{pos}}P_a$ ou $Z'>_{\mathrm{pos}}P_a$. Si $Z'<_{\mathrm{pos}}P_a$, on aura $Z'<_{\mathrm{pos}}\Psi$, donc $Z'\parallel\Psi$, on gagne encore.

% slanted
\margin{\uline{Lemme} (trivial) $X<_{\mathrm{pos}}\Psi\Longrightarrow\unsure{\operatorname{or.}X|\operatorname{or}_i X}<_{\mathrm{pos}}\Psi$}
